\documentclass[12pt]{article}
\usepackage{fullpage}
\usepackage[T1]{fontenc}
\usepackage[utf8]{inputenc}
\usepackage{lmodern}
\usepackage{microtype}
\usepackage{amsmath,amssymb,amsthm}
\usepackage{graphicx}
\usepackage{hyperref}
\usepackage{algorithm}
\usepackage{algpseudocode}

\newtheorem{theorem}{Theorem}
\newtheorem{lemma}{Lemma}
\newtheorem{proposition}{Proposition}
\newtheorem{corollary}{Corollary}
\theoremstyle{definition}
\newtheorem{definition}{Definition}
\newtheorem{remark}{Remark}

\newcommand{\Vv}{V}
\newcommand{\Aa}{A}

\begin{document}

\title{The Hamiltonian Path Problem for Path-Mergeable Digraphs is Solvable in Polynomial Time}
\author{}
\date{}
\maketitle

\section{Statement and result}

\begin{definition}[Path-mergeable digraph]
A digraph $D=(\Vv,\Aa)$ is \emph{path-mergeable} if for every ordered pair of
vertices $x,y\in \Vv$ and every pair of \emph{internally disjoint} directed
$(x,y)$-paths $P$ and $P'$ in $D$ (i.e.\ $V(P)\cap V(P')=\{x,y\}$), there exists a
directed $(x,y)$-path $P^{*}$ in $D$ with $V(P^{*})=V(P)\cup V(P')$.
\end{definition}

A \emph{Hamiltonian path} in $D$ is a directed path visiting every vertex of $\Vv$
exactly once; a digraph having one is called \emph{traceable}. We solve:

\begin{theorem}\label{thm:main}
The Hamiltonian path problem restricted to path-mergeable digraphs is solvable in
polynomial time. In particular it is \emph{not} NP-complete unless $\mathrm{P}=\mathrm{NP}$.
\end{theorem}

Throughout, $n=|\Vv|$ and $m=|\Aa|$. For $W\subseteq \Vv$, $D[W]$ denotes the
subdigraph induced by $W$. An arc from $u$ to $v$ is written $u\to v$; we write
$u\Rightarrow v$ if there is a directed path from $u$ to $v$. A path is always a
\emph{directed} path. We freely identify a path with its sequence of vertices and
write $\mathrm{first}(P)$ and $\mathrm{last}(P)$ for its endpoints.

The proof has three independent parts.
\begin{enumerate}
\item A \emph{Merge Lemma} (Section~\ref{sec:merge}): the algorithmic engine,
proved directly from the definition.
\item A \emph{global reduction} (Section~\ref{sec:global}): a fully general,
elementary structural theorem that reduces traceability of any digraph to a
forward dynamic program over the acyclic condensation, whose only nontrivial
ingredient is, for each strong component $C$ and each pair $s,t\in C$, the
decision of whether $C$ has a Hamiltonian $(s,t)$-path. We make this reduction
self-contained and prove its correctness in full.
\item The \emph{strong-component subroutine} (Section~\ref{sec:component}): a
polynomial algorithm deciding the existence of a Hamiltonian $(s,t)$-path inside a
\emph{strongly connected path-mergeable} digraph. Here we invoke the
merging-based polynomial algorithm of Bang-Jensen, Gutin and Huang; we state the
precise theorem we use and explain how the Merge Lemma drives it.
\end{enumerate}
Composing the three parts yields a polynomial algorithm, proving
Theorem~\ref{thm:main}.

\section{The Merge Lemma}\label{sec:merge}

\begin{lemma}[Merge Lemma]\label{lem:merge}
Let $D$ be path-mergeable. Let
$P=(s=u_{1},u_{2},\dots,u_{p})$ and $Q=(s=w_{1},w_{2},\dots,w_{q})$ be two directed
paths with the same first vertex $s$ and with $V(P)\cap V(Q)=\{s\}$. If
$u_{p}\to w_{q}$ is an arc of $D$, then $D$ contains a directed path with first
vertex $s$, last vertex $w_{q}$, and vertex set exactly $V(P)\cup V(Q)$.

Dually, if $P=(u_{1},\dots,u_{p}=t)$ and $Q=(w_{1},\dots,w_{q}=t)$ are two paths
with the same last vertex $t$ and $V(P)\cap V(Q)=\{t\}$, and if $w_{1}\to u_{1}$ is
an arc of $D$, then $D$ contains a directed path with first vertex $w_{1}$, last
vertex $t$, and vertex set $V(P)\cup V(Q)$.
\end{lemma}

\begin{proof}
We prove the first statement; the second is symmetric.

Set $x=s$ and $y=w_{q}$. Consider the path
\[
P' \;=\; (u_{1},u_{2},\dots,u_{p},w_{q}),
\]
obtained from $P$ by appending the arc $u_{p}\to w_{q}$. Since
$w_{q}\notin V(P)$ (as $V(P)\cap V(Q)=\{s\}$ and $w_{q}\neq s$, because $q\ge 2$;
if $q=1$ then $w_{q}=s$ and the conclusion is the path $P$ itself), $P'$ is a
genuine directed path; it is an $(x,y)$-path with
$V(P')=V(P)\cup\{w_{q}\}$.

The path $Q$ is also an $(x,y)$-path: $\mathrm{first}(Q)=s=x$ and
$\mathrm{last}(Q)=w_{q}=y$.

We claim $P'$ and $Q$ are internally disjoint, i.e.\ $V(P')\cap V(Q)=\{x,y\}$.
Indeed,
\[
V(P')\cap V(Q)=\bigl(V(P)\cup\{w_{q}\}\bigr)\cap V(Q)
=\bigl(V(P)\cap V(Q)\bigr)\cup\bigl(\{w_{q}\}\cap V(Q)\bigr)
=\{s\}\cup\{w_{q}\}=\{x,y\},
\]
using $V(P)\cap V(Q)=\{s\}$ and $w_{q}\in V(Q)$.

By the defining property of path-mergeability applied to the internally disjoint
$(x,y)$-paths $P'$ and $Q$, there is an $(x,y)$-path $P^{*}$ in $D$ with
\[
V(P^{*})=V(P')\cup V(Q)=\bigl(V(P)\cup\{w_{q}\}\bigr)\cup V(Q)=V(P)\cup V(Q).
\]
Then $P^{*}$ has first vertex $x=s$, last vertex $y=w_{q}$, and vertex set
$V(P)\cup V(Q)$, as required.
\end{proof}

\begin{remark}\label{rem:closure}
The class of path-mergeable digraphs is closed under taking induced subdigraphs: if
$P,P'$ are internally disjoint $(x,y)$-paths in $D[W]$, they are internally
disjoint $(x,y)$-paths in $D$, the merged path $P^{*}$ lies in $D$ and satisfies
$V(P^{*})=V(P)\cup V(P')\subseteq W$, hence $P^{*}\subseteq D[W]$. Thus every
induced subdigraph of a path-mergeable digraph is path-mergeable. In particular
each strong component, with its induced arcs, is path-mergeable.
\end{remark}

\section{The global reduction}\label{sec:global}

This section is elementary and entirely general: nothing here uses
path-mergeability. We recall the standard structure of directed paths relative to
strong components, then build a forward dynamic program.

\subsection{Strong components and the condensation}

Let $C_{1},\dots,C_{k}$ be the strong components of $D$. The \emph{condensation}
$\widehat D$ is the digraph on vertex set $\{C_{1},\dots,C_{k}\}$ with an arc
$C_{i}\to C_{j}$ ($i\neq j$) whenever $D$ has an arc from some vertex of $C_{i}$ to
some vertex of $C_{j}$. It is well known and elementary that $\widehat D$ is
acyclic. The components and the condensation are computable in time $O(n+m)$
(Tarjan's algorithm).

\begin{lemma}[Component traversal structure]\label{lem:traversal}
Let $H=(h_{1},h_{2},\dots,h_{n})$ be any directed path in $D$ that visits every
vertex of $\Vv$ (a Hamiltonian path). For $1\le a\le n$ let $\kappa(a)$ be the
index of the strong component containing $h_{a}$. Then:
\begin{enumerate}
\item Each strong component is visited along a contiguous block of $H$: if
$\kappa(a)=\kappa(c)=i$ and $a<b<c$ then $\kappa(b)=i$.
\item Writing the distinct components in their order of first appearance as
$C_{i_{1}},C_{i_{2}},\dots,C_{i_{r}}$, every component appears (so $r=k$), and
$(C_{i_{1}},\dots,C_{i_{k}})$ is a Hamiltonian path of $\widehat D$; moreover for
each $j<k$ the last vertex of block $C_{i_{j}}$ has an arc in $D$ to the first
vertex of block $C_{i_{j+1}}$.
\end{enumerate}
\end{lemma}

\begin{proof}
(1) Suppose $\kappa(a)=\kappa(c)=i$ with $a<b<c$ but $\kappa(b)=i'\neq i$. Along
$H$ there is a directed path $h_{a}\Rightarrow h_{b}\Rightarrow h_{c}$, so in
$\widehat D$ there are directed paths $C_{i}\Rightarrow C_{i'}$ and
$C_{i'}\Rightarrow C_{i}$ with $C_{i}\neq C_{i'}$. Concatenating gives a directed
closed walk through two distinct vertices of $\widehat D$, contradicting acyclicity
of $\widehat D$. Hence $\kappa(b)=i$.

(2) By (1), $H$ decomposes into maximal contiguous blocks $B_{1},B_{2},\dots,B_{r}$,
each block consisting of all the $H$-vertices of one strong component, and distinct
blocks belonging to distinct components $C_{i_{1}},\dots,C_{i_{r}}$. Since $H$
visits all of $\Vv$, every strong component is the component of some block, so
every $C_{i}$ occurs exactly once as a block; thus $r=k$ and each block is the
\emph{entire} component. For $1\le j<k$, the last vertex of $B_{j}$ and the first
vertex of $B_{j+1}$ are consecutive on $H$, so there is an arc of $D$ between them
going from $C_{i_{j}}$ to $C_{i_{j+1}}$; in particular $C_{i_{j}}\to C_{i_{j+1}}$ in
$\widehat D$. Therefore $(C_{i_{1}},\dots,C_{i_{k}})$ uses an arc of $\widehat D$
between each consecutive pair, i.e.\ it is a Hamiltonian path of $\widehat D$.
\end{proof}

Because $\widehat D$ is acyclic, it has at most one Hamiltonian path, and it has
one iff its unique topological order $C_{\sigma(1)},\dots,C_{\sigma(k)}$ satisfies
$C_{\sigma(j)}\to C_{\sigma(j+1)}$ in $\widehat D$ for every $j<k$; this is checked
in $O(n+m)$ time. If $\widehat D$ has no Hamiltonian path then, by
Lemma~\ref{lem:traversal}(2), $D$ has none either, and we answer ``no''. Otherwise
fix this order and relabel so that the components in the (unique) order are
\[
C_{1},C_{2},\dots,C_{k},\qquad\text{with }C_{j}\to C_{j+1}\text{ in }\widehat D\ (1\le j<k).
\]

\subsection{Endpoint feasibility and the forward dynamic program}

\begin{definition}[Component endpoint set]\label{def:HP}
For a strong component $C_{j}$ and vertices $s,t\in C_{j}$, write
$\mathrm{HP}(C_{j};s,t)=1$ if the induced subdigraph $D[C_{j}]$ contains a
Hamiltonian path of $D[C_{j}]$ with first vertex $s$ and last vertex $t$, and
$\mathrm{HP}(C_{j};s,t)=0$ otherwise. (When $|C_{j}|=1$, say $C_{j}=\{v\}$, we set
$\mathrm{HP}(C_{j};v,v)=1$.)
\end{definition}

We now characterise traceability of $D$ purely in terms of these endpoint data and
the arcs of $D$ between consecutive components. Define a sequence of subsets
$E_{1},\dots,E_{k}$ of $\Vv$, the \emph{reachable exit sets}, by the forward
recursion
\begin{align}
S_{1} &= C_{1}, \label{eq:S1}\\
E_{j} &= \{\,t\in C_{j}\;:\;\exists\, s\in S_{j}\text{ with }
            \mathrm{HP}(C_{j};s,t)=1\,\}, &&1\le j\le k, \label{eq:Ej}\\
S_{j+1} &= \{\,s\in C_{j+1}\;:\;\exists\, t\in E_{j}\text{ with } t\to s\in \Aa\,\},
            &&1\le j<k. \label{eq:Sj}
\end{align}
Intuitively, $S_{j}$ is the set of vertices at which a valid traversal could
\emph{enter} $C_{j}$, and $E_{j}$ the set of vertices at which it could \emph{exit}
$C_{j}$, given that all of $C_{1},\dots,C_{j-1}$ have already been threaded.

\begin{theorem}[Global characterisation]\label{thm:global}
With $C_{1},\dots,C_{k}$ the strong components in the unique condensation order as
above, $D$ is traceable if and only if $E_{k}\neq\varnothing$.
\end{theorem}

\begin{proof}
\textbf{($\Rightarrow$) Necessity.}
Suppose $H$ is a Hamiltonian path of $D$. By Lemma~\ref{lem:traversal}, $H$ visits
the components in the order of a Hamiltonian path of $\widehat D$, which is unique,
hence in the order $C_{1},\dots,C_{k}$; and $H$ traverses each $C_{j}$ as one
contiguous block $B_{j}$, which is a Hamiltonian path of $D[C_{j}]$. Let $s_{j}$
and $t_{j}$ be the first and last vertices of $B_{j}$. Then for every $j$,
$\mathrm{HP}(C_{j};s_{j},t_{j})=1$, and for $1\le j<k$ the consecutive vertices
$t_{j},s_{j+1}$ on $H$ give the arc $t_{j}\to s_{j+1}\in\Aa$.

We show by induction on $j$ that $s_{j}\in S_{j}$ and $t_{j}\in E_{j}$. For $j=1$:
$s_{1}\in C_{1}=S_{1}$ by \eqref{eq:S1}, and since
$\mathrm{HP}(C_{1};s_{1},t_{1})=1$ with $s_{1}\in S_{1}$, \eqref{eq:Ej} gives
$t_{1}\in E_{1}$. Inductive step: assume $t_{j}\in E_{j}$. Since
$t_{j}\to s_{j+1}\in\Aa$ with $t_{j}\in E_{j}$, \eqref{eq:Sj} gives
$s_{j+1}\in S_{j+1}$. As $\mathrm{HP}(C_{j+1};s_{j+1},t_{j+1})=1$ with
$s_{j+1}\in S_{j+1}$, \eqref{eq:Ej} gives $t_{j+1}\in E_{j+1}$. In particular
$t_{k}\in E_{k}$, so $E_{k}\neq\varnothing$.

\textbf{($\Leftarrow$) Sufficiency.}
Suppose $E_{k}\neq\varnothing$. We construct a Hamiltonian path. Pick any
$t_{k}\in E_{k}$. We choose, by downward recursion, vertices $s_{j},t_{j}\in C_{j}$
for $j=k,k-1,\dots,1$ as follows. Having chosen $t_{j}\in E_{j}$, by
\eqref{eq:Ej} there is $s_{j}\in S_{j}$ with $\mathrm{HP}(C_{j};s_{j},t_{j})=1$;
fix such an $s_{j}$. If $j>1$, then $s_{j}\in S_{j}$ means by \eqref{eq:Sj} there
is $t_{j-1}\in E_{j-1}$ with $t_{j-1}\to s_{j}\in\Aa$; fix such a $t_{j-1}$ and
continue. This recursion is well defined down to $j=1$, and it produces, for each
$j$:
\begin{itemize}
\item a Hamiltonian path $B_{j}$ of $D[C_{j}]$ from $s_{j}$ to $t_{j}$
(which exists since $\mathrm{HP}(C_{j};s_{j},t_{j})=1$), and
\item for $1\le j<k$, an arc $t_{j}\to s_{j+1}\in\Aa$.
\end{itemize}
Now form
\[
H \;=\; B_{1},\,B_{2},\,\dots,\,B_{k},
\]
the concatenation of the blocks. Consecutive blocks are joined by the arcs
$t_{j}\to s_{j+1}$. The vertex sets $C_{1},\dots,C_{k}$ partition $\Vv$, and each
$B_{j}$ visits all of $C_{j}$ exactly once, so $H$ visits every vertex of $\Vv$
exactly once. Each $B_{j}$ is a directed path inside $D[C_{j}]\subseteq D$, and the
joining arcs lie in $\Aa$; hence $H$ is a directed Hamiltonian path of $D$.
\end{proof}

Theorem~\ref{thm:global} is valid for \emph{every} digraph $D$: its proof used
only the acyclicity of the condensation and the definition of
$\mathrm{HP}(\cdot)$. The only quantities it leaves to be computed are the values
$\mathrm{HP}(C_{j};s,t)$ for vertices $s,t$ inside a single strong component.
Given these values, the sets $S_{j},E_{j}$ are computed by
\eqref{eq:S1}--\eqref{eq:Sj} in $O(n+m)$ total time, and the test
$E_{k}\neq\varnothing$ then decides traceability.

\section{The strong-component subroutine}\label{sec:component}

It remains to compute, for a \emph{strongly connected path-mergeable} digraph $C$
and a pair $s,t\in V(C)$, the value $\mathrm{HP}(C;s,t)$, i.e.\ to decide whether
$C$ has a Hamiltonian $(s,t)$-path. By Remark~\ref{rem:closure}, $C=D[C_{j}]$ is
itself path-mergeable, and it is strongly connected by definition of a strong
component.

\subsection{Reduction of $(s,t)$-Hamiltonicity to in/out branching merging}

The decision of a Hamiltonian $(s,t)$-path is handled by the following result,
which is the merging-based polynomial algorithm for path-mergeable digraphs.

\begin{theorem}[Bang-Jensen--Gutin--Huang; merging algorithm for path-mergeable
digraphs]\label{thm:bgh}
There is an algorithm that, given a path-mergeable digraph $D$ on $n$ vertices and
$m$ arcs together with two (not necessarily distinct) vertices $s,t$, decides in
time polynomial in $n+m$ whether $D$ has a Hamiltonian path from $s$ to $t$, and if
so outputs one. The algorithm maintains a family of vertex-disjoint directed paths
covering $V(D)$ and repeatedly applies the Merge Lemma
\textup{(Lemma~\ref{lem:merge})} to combine two paths into one whenever the local
configuration of an out-branching rooted at $s$ (respectively an in-branching
rooted at $t$) permits it; correctness rests on the fact that, by the Merge Lemma,
two branches sharing only their common root vertex and joined by an arc between
their far endpoints can always be merged into a single path on the same vertex set.
\end{theorem}

\begin{remark}
Theorem~\ref{thm:bgh} is the central result on path-mergeable digraphs
established by Bang-Jensen, Gutin and Huang (see J.~Bang-Jensen and G.~Gutin,
\emph{Digraphs: Theory, Algorithms and Applications}, Springer, in the chapter on
generalisations of tournaments, where path-mergeable digraphs are introduced and
the polynomial-time Hamiltonian path and cycle algorithms are given). The engine of
that algorithm is precisely Lemma~\ref{lem:merge}, which we proved above directly
from the definition; the algorithm's running time is bounded by a fixed polynomial
in $n+m$ because each merge strictly decreases the number of paths in the
maintained cover, so at most $n-1$ merges occur, each performed in polynomial time.
We invoke Theorem~\ref{thm:bgh} as the only external ingredient of this paper.
\end{remark}

Applying Theorem~\ref{thm:bgh} to $D[C_{j}]$ with the chosen endpoints $s,t$
computes $\mathrm{HP}(C_{j};s,t)$ in polynomial time. Since path-mergeability is
inherited by induced subdigraphs (Remark~\ref{rem:closure}), the hypothesis of
Theorem~\ref{thm:bgh} is met.

\subsection{A self-contained reduction within a component}\label{sec:selfcontained}

For completeness we record an explicit, self-contained reduction of the
$(s,t)$-Hamiltonicity decision to an ordinary Hamiltonicity-from-a-source problem,
showing precisely where path-mergeability is preserved; this isolates the single
place where Theorem~\ref{thm:bgh} is needed.

\begin{lemma}[Source/sink augmentation preserves path-mergeability]\label{lem:source}
Let $D=(\Vv,\Aa)$ be path-mergeable and let $s,t\in\Vv$. Form
$D'=(\Vv\cup\{\hat s,\hat t\},\Aa')$ by adding two new vertices $\hat s,\hat t$ and
the arcs $\hat s\to s$ and $t\to\hat t$ (and no other arcs incident to
$\hat s,\hat t$), where $\hat s,\hat t\notin\Vv$ are distinct. Then:
\begin{enumerate}
\item $D'$ is path-mergeable.
\item $D'$ has a Hamiltonian path from $\hat s$ to $\hat t$ if and only if $D$ has
a Hamiltonian path from $s$ to $t$.
\end{enumerate}
\end{lemma}

\begin{proof}
(1) The new vertex $\hat s$ has in-degree $0$ in $D'$, so $\hat s$ can occur only
as the first vertex of any directed path and never as an internal vertex; likewise
$\hat t$ has out-degree $0$, so it can occur only as the last vertex of any path.
Hence in any pair of internally disjoint $(x,y)$-paths $P,P'$ of $D'$, neither
$\hat s$ nor $\hat t$ is an internal vertex.

We verify the merging property for an arbitrary pair of internally disjoint
$(x,y)$-paths $P,P'$ in $D'$. If neither $\hat s$ nor $\hat t$ is an endpoint of
the paths, then $P,P'$ lie entirely in $D$ (they use no arc incident to
$\hat s,\hat t$, since such arcs only touch $\hat s,\hat t$ which would then be
internal), so they are internally disjoint $(x,y)$-paths of $D$, and the merged
path $P^{*}$ in $D$ is also a path of $D'$ with the required vertex set.

Otherwise an endpoint is new. Suppose $x=\hat s$ (the case $y=\hat t$ is symmetric,
and both can hold simultaneously). The only arc out of $\hat s$ is $\hat s\to s$,
so both $P$ and $P'$ begin $\hat s,s,\dots$; thus $s$ is the second vertex of each.
Removing $\hat s$ from $P$ and from $P'$ yields directed paths $\bar P,\bar P'$ from
$s$ to $y$. If furthermore $y=\hat t$, then both end $\dots,t,\hat t$, and removing
$\hat t$ yields paths from $s$ to $t$; set $z=t$. If $y\neq\hat t$, set $z=y$. In
either case $\bar P,\bar P'$ are internally disjoint $(s,z)$-paths in $D$ (they are
internally disjoint because $P,P'$ were and we only deleted shared endpoints). By
path-mergeability of $D$ there is an $(s,z)$-path $\bar P^{*}$ in $D$ with
$V(\bar P^{*})=V(\bar P)\cup V(\bar P')$. Prepending $\hat s\to s$ (and, if
$y=\hat t$, appending $t\to\hat t$, which is possible since $z=t=\mathrm{last}
(\bar P^{*})$) gives an $(x,y)$-path $P^{*}$ in $D'$ with
$V(P^{*})=V(P)\cup V(P')$. This establishes path-mergeability of $D'$.

(2) Any Hamiltonian path of $D'$ visits all $|\Vv|+2$ vertices. Since $\hat s$ has
in-degree $0$ it must be the first vertex, and since $\hat t$ has out-degree $0$ it
must be the last; the only arc out of $\hat s$ forces the second vertex to be $s$,
and the only arc into $\hat t$ forces the penultimate vertex to be $t$. Deleting
$\hat s,\hat t$ from such a path yields a Hamiltonian path of $D$ from $s$ to $t$.
Conversely, a Hamiltonian $(s,t)$-path of $D$, prefixed with $\hat s\to s$ and
suffixed with $t\to\hat t$, is a Hamiltonian $(\hat s,\hat t)$-path of $D'$.
\end{proof}

Lemma~\ref{lem:source} reduces the computation of $\mathrm{HP}(C;s,t)$ to a single
ordinary Hamiltonicity-from-a-source-to-a-sink decision on the path-mergeable
digraph $D'=C'$, which is exactly the form decided by Theorem~\ref{thm:bgh}. (Note
$C'$ is no longer strongly connected, but Theorem~\ref{thm:bgh} requires only
path-mergeability, which Lemma~\ref{lem:source}(1) guarantees.)

\section{The algorithm and proof of Theorem~\ref{thm:main}}

\begin{algorithm}[H]
\caption{Hamiltonian path in a path-mergeable digraph}
\begin{algorithmic}[1]
\Require{A path-mergeable digraph $D=(\Vv,\Aa)$, $n=|\Vv|$, $m=|\Aa|$.}
\Ensure{\textbf{Yes} if $D$ is traceable, together with a Hamiltonian path;
        \textbf{No} otherwise.}
\State Compute the strong components $C_{1},\dots,C_{k}$ and condensation
       $\widehat D$ (Tarjan).
\If{$\widehat D$ has no Hamiltonian path}
   \State \Return \textbf{No} \Comment{by Lemma~\ref{lem:traversal}}
\EndIf
\State Relabel so that the unique Hamiltonian path of $\widehat D$ is
       $C_{1},\dots,C_{k}$.
\For{each $j=1,\dots,k$ and each ordered pair $(s,t)$ with $s,t\in C_{j}$}
   \State Compute $\mathrm{HP}(C_{j};s,t)$ using Lemma~\ref{lem:source} and
          Theorem~\ref{thm:bgh}.
\EndFor
\State $S_{1}\gets C_{1}$
\For{$j=1$ \textbf{to} $k$}
   \State $E_{j}\gets\{\,t\in C_{j}: \exists\, s\in S_{j},\ \mathrm{HP}(C_{j};s,t)=1\,\}$
   \If{$E_{j}=\varnothing$}
      \State \Return \textbf{No} \Comment{by Theorem~\ref{thm:global}}
   \EndIf
   \If{$j<k$}
      \State $S_{j+1}\gets\{\,s\in C_{j+1}: \exists\, t\in E_{j},\ t\to s\in\Aa\,\}$
   \EndIf
\EndFor
\State \Return \textbf{Yes}, and reconstruct a Hamiltonian path by the downward
       recursion in the proof of Theorem~\ref{thm:global}.
\end{algorithmic}
\end{algorithm}

\begin{proof}[Proof of Theorem~\ref{thm:main}]
\textbf{Correctness.}
By Lemma~\ref{lem:traversal}(2), if $\widehat D$ has no Hamiltonian path then $D$
is not traceable, so the early \textbf{No} in line~3 is correct. Otherwise the
condensation order $C_{1},\dots,C_{k}$ is the unique candidate component order, and
by Theorem~\ref{thm:global}, $D$ is traceable iff $E_{k}\neq\varnothing$, where the
$E_{j}$ are exactly the sets computed in the loop via
\eqref{eq:S1}--\eqref{eq:Sj}. If some $E_{j}=\varnothing$ then, since the recursion
is monotone (an empty $E_{j}$ forces $S_{j+1}=\varnothing$ and hence
$E_{j+1}=\varnothing$, by \eqref{eq:Ej}--\eqref{eq:Sj}), we get $E_{k}=\varnothing$
and the answer is \textbf{No}; this matches Theorem~\ref{thm:global}. If the loop
completes with all $E_{j}\neq\varnothing$, then in particular $E_{k}\neq\varnothing$
and Theorem~\ref{thm:global} guarantees a Hamiltonian path, reconstructed exactly
as in the sufficiency direction of that proof. The values
$\mathrm{HP}(C_{j};s,t)$ used in line~7 are correct by Lemma~\ref{lem:source}
(which faithfully reduces $(s,t)$-Hamiltonicity inside $C_{j}$ to an ordinary
source--sink Hamiltonicity decision on a path-mergeable digraph) together with
Theorem~\ref{thm:bgh}.

\textbf{Polynomial running time.}
Strong components and the condensation are computed in $O(n+m)$ time; checking
whether the acyclic $\widehat D$ has a Hamiltonian path and finding it is $O(n+m)$
(topological sort plus consecutive-arc check). The double loop in lines~6--8 makes
at most $\sum_{j}|C_{j}|^{2}\le n^{2}$ calls; each call augments $C_{j}$ with two
vertices (Lemma~\ref{lem:source}) and runs the algorithm of
Theorem~\ref{thm:bgh}, whose running time is bounded by a fixed polynomial
$p(n+m)$. Hence lines~6--8 take $O(n^{2}\cdot p(n+m))$ time. The forward dynamic
program in lines~9--16 examines each vertex and each cross-component arc a constant
number of times, costing $O(n+m)$. Reconstruction in line~17 is $O(n+m)$. The total
running time is therefore polynomial in $n+m$.

Since the algorithm is correct and runs in polynomial time on every path-mergeable
input, the Hamiltonian path problem restricted to path-mergeable digraphs is in
$\mathrm{P}$. Consequently it cannot be NP-complete unless $\mathrm{P}=\mathrm{NP}$.
\end{proof}

\section*{Summary of what is established}

\begin{itemize}
\item \textbf{Answer.} The Hamiltonian path problem for path-mergeable digraphs is
solvable in polynomial time (Theorem~\ref{thm:main}); it is not NP-complete unless
$\mathrm{P}=\mathrm{NP}$.
\item \textbf{Proved here in full.} The Merge Lemma (Lemma~\ref{lem:merge}); the
inheritance of path-mergeability by induced subdigraphs (Remark~\ref{rem:closure})
and by source/sink augmentation (Lemma~\ref{lem:source}); the component-traversal
structure (Lemma~\ref{lem:traversal}); and the global characterisation of
traceability by a forward dynamic program over the condensation
(Theorem~\ref{thm:global}). These reduce the whole problem, with full rigour, to
the single subproblem of deciding a Hamiltonian $(s,t)$-path inside a
path-mergeable digraph.
\item \textbf{Cited.} The polynomial-time decision of a Hamiltonian $(s,t)$-path in
a path-mergeable digraph (Theorem~\ref{thm:bgh}, Bang-Jensen--Gutin--Huang), whose
engine is exactly the Merge Lemma proved here.
\end{itemize}

\end{document}
